\documentclass{standalone}
\usepackage{circuitikz}
\usetikzlibrary{calc}
\definecolor{circuitnotemark}{HTML}{FE00FE}
\begin{document}
\begin{circuitikz}[american, line width=0.8pt]
\ctikzset{bipoles/length=1.2cm}
\foreach \x in {0,2,4,6,8,10,12} {\foreach \y in {0} {\fill[gray, opacity=0.35] (\x,\y) circle (1.1pt);}}
\node[gray, font=\scriptsize] at (0,0.84) {1};
\node[gray, font=\scriptsize] at (2,0.84) {2};
\node[gray, font=\scriptsize] at (4,0.84) {3};
\node[gray, font=\scriptsize] at (6,0.84) {4};
\node[gray, font=\scriptsize] at (8,0.84) {5};
\node[gray, font=\scriptsize] at (10,0.84) {6};
\node[gray, font=\scriptsize] at (12,0.84) {7};
\node[gray, font=\scriptsize, anchor=east] at (-0.84,0) {a};
\coordinate (a1) at (0,0);
\coordinate (a3) at (4,0);
\coordinate (a5) at (8,0);
\coordinate (a7) at (12,0);
\draw (a1) to[R, l_=$R_{1}$, a^=$10\,\mathrm{k}\Omega$] (a3); % line 3
\draw (a5) to[L, l_=$L_{1}$, a^=$10\,\mathrm{m}\mathrm{H}$] (a7); % line 4
\path (14,-3.641) rectangle (19.103,0.395); % line 6
\node[anchor=west, inner sep=0, circuitnotemark, font=\footnotesize] at (14,0) {X}; % line 6
\node[anchor=west, inner sep=0, circuitnotemark, font=\footnotesize] at (14,-0.325) {X}; % line 6
\node[anchor=west, inner sep=0, circuitnotemark, font=\footnotesize] at (14,-0.649) {X}; % line 6
\node[anchor=west, inner sep=0, circuitnotemark, font=\footnotesize] at (14,-0.974) {X}; % line 6
\node[anchor=west, inner sep=0, circuitnotemark, font=\footnotesize] at (14,-1.298) {X}; % line 6
\node[anchor=west, inner sep=0, circuitnotemark, font=\footnotesize] at (14,-1.623) {X}; % line 6
\node[anchor=west, inner sep=0, circuitnotemark, font=\footnotesize] at (14,-1.947) {X}; % line 6
\node[anchor=west, inner sep=0, circuitnotemark, font=\footnotesize] at (14,-2.272) {X}; % line 6
\node[anchor=west, inner sep=0, circuitnotemark, font=\footnotesize] at (14,-2.596) {X}; % line 6
\node[anchor=west, inner sep=0, circuitnotemark, font=\footnotesize] at (14,-2.921) {X}; % line 6
\node[anchor=west, inner sep=0, circuitnotemark, font=\footnotesize] at (14,-3.246) {X}; % line 6
\node[anchor=south west, inner sep=2pt, circuitnotemark, font=\large\bfseries] (circuittitle) at (current bounding box.north west) {X};
\path (circuittitle.south west) rectangle ++(7.493,0.593);
\end{circuitikz}
\end{document}
